\uuid{ClIp}
\titre{Logique}
\auteur{bodin ; croizat ; sacre}
\organisation{exo7}
\licence{CC BY-NC-SA 4.0}
\datecreate{2022-03-08}
\niveau{L1}
\module{Algèbre}
\chapitre{Logique, ensemble, raisonnement}
\sousChapitre{Logique}
\difficulte{2}
\theme{Logique}
\exo7qcmid{hilisit.logique.009}

\contenu{
\begin{question}
\qclassification{100.01}
On considère l'implication :
$$\text{"tu prépares un repas } \implies \text{ je viens chez toi".}$$
Quelles sont les affirmations vraies ?
\begin{answers}
	\bad{L'implication réciproque est "je viens chez toi $\implies$ tu ne prépares pas de repas".}
	\good{La contraposée est "je ne viens pas chez toi $\implies$  tu ne prépares pas de repas".}
	\bad{Si l'implication est vraie alors l'implication réciproque l'est aussi.}
	\good{Si l'implication est vraie alors sa contraposée l'est aussi.}
\end{answers}
\begin{explanations}
L'implication réciproque est "je viens chez toi $\implies$ tu prépares un repas". \\
La contraposée est "je ne viens pas chez toi $\implies$ tu ne prépares pas de repas".\\
Une implication et sa contraposée sont des propositions équivalentes (donc vraies en même temps, et fausses en même temps). \\
\end{explanations}
\end{question}
}
